\paragraph{Misuse and limitations.}
\notest{} must not be read as ``no effect'' or as evidence that surviving cells
are behaving normally. A separate within-cell analysis over three cohorts and
42 patients also remained unresolved because its comparability checks returned
\emph{undecided} (Appendix~\ref{app:guards}). The evidence consists of one empirical self-audit and a constructed
external-control example; neither validates a patient world model or a
replacement reporting rule. Calibration covers one gene,
one labelling axis and one weighting on two small cohorts. It challenges this
rule's justification, not cell-count rules in general. The fixed-pool analysis estimates coverage conditional on empirical cell
distributions; uncertainty over new patients remains unmeasured. The draw-pool choice was
underspecified, repeated simulations do not add independent patients, and
the leave-one-patient-out sensitivity remains conditional on these cohorts;
patient-level population uncertainty in the candidate cutpoints is unquantified.
The absence-defined labelling axis remains sensitive to sequencing depth
(Appendix~\ref{app:limitations}). The external search found no affirmative case, so the documented failure
remains a self-audit rather than evidence about its prevalence. The external-control
example uses a synthetic registry and an illustrative decision threshold;
it does not establish clinical benefit or support a treatment decision.

\paragraph{Data and impact.}
The single-cell cohorts are public, de-identified and previously published.
The release policy covers derived summaries and code, rather than cell-level
matrices. A system reporting group-specific quantities should distinguish an
unavailable comparison from an estimated null and should evaluate the coverage
and discrimination of any abstention rule. Publishing failed calibration
attempts helps prevent unsupported thresholds from being treated as validated.
